% Where the latency actually is in closed-loop EEG-driven music
%
% Compiles with pdflatex (no unicode, no exotic packages):
%     pdflatex arxiv_short.tex && pdflatex arxiv_short.tex
% Two passes so the hyperref/table references settle.
%
% Source of truth for the numbers is docs/arxiv_short.md and, behind it,
% `python scripts/verify_claims.py`. If a number here disagrees with that, the
% script is right.

\documentclass[11pt,a4paper]{article}

\usepackage[T1]{fontenc}
\usepackage[utf8]{inputenc}
\usepackage[margin=1in]{geometry}
\usepackage{booktabs}
\usepackage{amsmath}
\usepackage{microtype}
\usepackage[hidelinks]{hyperref}

\title{Where the latency actually is in closed-loop EEG-driven music\\
       \large A measured architecture, and four statistical properties of
       closed-loop EEG data}
\author{Kirthan Kulkarni}
\date{}

\begin{document}
\maketitle

\begin{abstract}
Closed-loop systems that drive music generation from EEG are an active area, and a working
one: several have been built and described. What is not reported is how long they take.
Across the systems we could resolve individually, none publishes a measured, decomposed
end-to-end latency budget; one reports its analysis window and update rate and reasons about
delay qualitatively.

We build such a system on consumer hardware---a four-channel Muse headband, a frontal
$\log(\beta/\alpha)$ index $z$-scored against the participant's own baseline, a nine-rung
prompt ladder, and a resident library of pre-rendered audio---and measure it. Generation with
MusicGen runs $1.05$--$6.27\times$ slower than realtime across every machine, backend and
precision we tested, including a datacenter T4, so a precomputed library is not a compromise
but the correct architecture for a finite prompt space. End-to-end worst case is $6.5$~s, of
which $84.6\%$ is the analysis path, not the model.

That decomposition turns out to be configuration-dependent rather than architectural.
Against labelled eyes-open/closed ground truth, the deployed estimator is dominated by 8 of
its 9 alternatives on both detection latency and information per minute simultaneously. The
architecture's floor is $0.189$~s of analysis path, and at that floor information rate is
$1.83\times$ the deployed value---the fastest configuration is also the more informative one.
With a $0.25$~s crossfade the end-to-end figure is $0.439$~s, inside the regime where
published alpha-neurofeedback studies still find an effect and where our deployed
configuration is $13\times$ too slow.

We also report four statistical properties of closed-loop EEG data that affect any study of
this kind, three of which cost us designs before we measured them: autocorrelation inflates
apparent evidence $6.4\times$; continuous audio-neural coupling loses power precisely when
the intervention succeeds; event-locked responses are confounded by their own trigger to a
degree we calibrate against synthetic ground truth; and dichotomised outcomes cost roughly
$3\times$ in detectable effect.

\textbf{This is a systems and methods paper. It makes no claim that the intervention helps
anyone.} Two self-administered unblinded sessions exist; zero adaptive/sham pairs; zero
participants.
\end{abstract}

\section{What is missing from a crowded field}

Building an EEG-to-music loop is not a contribution. MindMelody~\cite{mindmelody} pairs a
transformer-GNN affect encoder with a retrieval-augmented planner over MusicGen-medium; Ran
et al.~\cite{mindtomusic} drive real-time emotional music generation from EEG; Ehrlich et
al.~\cite{ehrlich} close the loop with synthesised affective music; Velasquez-Pena et
al.~\cite{neurophone} do it on a consumer Muse S into Ableton Live; Monroy-D'Croz et
al.~\cite{minimalist} build a minimalist BCMI on prefrontal EEG over Lab Streaming Layer.
Dash and Agres~\cite{dashagres} survey the field.

These systems decode affect better than ours. We use a single hand-built index on two
electrodes; MindMelody reports $76.8\%$ valence and $72.4\%$ arousal cross-subject on DEAP.
Any claim of novelty in system capability is unavailable to us, and we make none.

What they do not report is timing. MindMelody describes itself as closed-loop and real-time
and gives no latency figure, no inference time and no hardware specification.
Monroy-D'Croz et al.---architecturally the closest published system to ours---reports none
either, and ran no control condition.

One paper is a partial counterexample and we state it rather than leave it for a reviewer.
Ehrlich et al.\ report a $4$~s analysis window at $87.5\%$ overlap giving a $0.5$~s update
rate, say that zero-phase filtering was chosen to avoid delay, and describe tuning a feedback
parameter against perceived latency. They convert none of this into a signal-to-audio figure.
So the defensible claim is not that this literature ignores latency; it is that \emph{no
system in this group reports a measured, decomposed end-to-end latency budget}.

The adjacent fields treat it as mandatory. Gesture2Music~\cite{gesture2music} reports $30$~ms
inference against a stated $100$~ms interaction threshold. Neurofeedback work treats delay as
the independent variable: Belinskaia et al.~\cite{belinskaia} trained parietal alpha with
feedback at $0$~ms, $+250$~ms and $+500$~ms and found sustained post-training change
negatively correlated with latency and absent at $500$~ms; Asai et al.~\cite{asai} found an
effect on EEG microstates only at zero delay, with one second enough to abolish it and
participants unable to detect the delay that did so.

Our question is therefore not whether such a system can be built. It is: when one is measured
rather than described, what is the number, and what follows from it?

\section{The system}

\paragraph{Sensing.} A Muse 2 streams four channels at $256$~Hz over LSL. Frontal AF7/AF8 are
averaged; each window is detrended, band-passed $1$--$45$~Hz, notched at $60$~Hz, and
Welch-transformed. The index is $\log(\beta/\alpha)$, smoothed with a one-pole filter
($\tau = 3$~s) and $z$-scored against the participant's own $120$~s eyes-open baseline.
Windows exceeding a $350~\mu$V peak-to-peak threshold are rejected and do not update the
smoother.

\paragraph{Control.} A nine-rung energy ladder spans arousal at $0.5$~SD per rung. The
controller applies the iso-principle~\cite{altshuler}: it does not jump to the target state
but moves one rung from the participant's current rung toward it. Output is therefore a pure
function of $(z, \mathrm{target})$, with a finite range of exactly nine strings.

\paragraph{Audio.} Each rung has $32$ pre-rendered $8$~s segments ($288$ total, $38.4$
minutes). A prompt change abandons the current segment mid-playback and equal-power
crossfades into one drawn from the new rung.

Two properties of this design carry the paper.

\emph{The prompt space is finite and small, so generation is unnecessary rather than merely
slow.} A library covering nine strings covers the controller's entire output space exactly.
Even given an infinitely fast generator, re-synthesising one of nine prompts is strictly
worse than selecting a pre-rendered variant. This argument survives arbitrarily fast models;
the latency argument below does not.

\emph{The library responds in one crossfade, not one segment.} Streaming commits: once a
segment is queued it plays to completion, so a prompt change waits up to
$\mathrm{queue\_depth} \times \mathrm{segment\_seconds}$ regardless of GPU speed. Worst-case
audio latency drops from $8.0$~s to $1.0$~s---$8\times$, and it is the architecture, not the
hardware, that buys it.

\section{The latency budget, measured}

\subsection{Generation is slower than realtime everywhere we looked}

Nine benchmark runs across two GPU tiers, two backends and three precision modes. Every one
of $44$ configuration cells lands between $1.05\times$ and $6.27\times$ slower than realtime.
The best case is a datacenter T4.

The bound is the sequential decode loop, not arithmetic, established by giving the workload
$8\times$ the arithmetic and watching it get slower: on the T4, fp32 beats fp16-half by a
factor of $0.952$, and fp32 $<$ fp16-half $<$ fp16 held in 6 of 6 cells. Autocast inserts a
cast at every operation, which amortises over large batched matmuls in training and is pure
overhead for batch-1 autoregressive decoding.

Between-run variance is a property of the machine, not the measurement: up to $1.96\times$ on
a thermally limited laptop and $1.18\times$ on the T4 for the same configuration. Report
ranges, not point estimates---a single run of this probe is not a measurement on a laptop GPU.

\subsection{The dominant term is the analysis path}

End-to-end worst case with the library engine is $6.5$~s: $5.5$~s of analysis plus one
$1.0$~s crossfade. $84.6\%$ is analysis. Work aimed at faster generation optimises the
smaller term.

The analysis budget is structural, not implementation slop: a $4$~s window contributes
$2.0$~s of centroid delay, a $1$~s hop contributes $0.5$~s of quantisation, and a
$\tau = 3$~s smoother contributes $3.0$~s of group delay. Predicted $5.5$~s; measured
$5.665$~s against labelled ground truth.

\subsection{The deployed configuration is dominated}

Latency alone would favour a wire and discriminability alone would favour smoothing
everything to a constant, so we score estimators on $d \times \sqrt{\text{independent
observations per minute}}$---the $t$-statistic per root-minute, which is what determines how
precisely a fixed-duration session resolves a state difference.

\begin{table}[h]
\centering
\begin{tabular}{lrrrr}
\toprule
estimator & detect & $d$ & ind/min & info/min \\
\midrule
deployed ($4$~s window, $1$~s hop, $\tau = 3$) & $5.67$~s & $1.99$ & $1.2$ & $2.14$ \\
$2$~s window, $0.5$~s hop, $\tau = 0.5$ & $1.68$~s & $1.14$ & $13.6$ & $4.20$ \\
streaming, order 4, $\tau = 0.25$ & $0.17$~s & $0.70$ & $30.9$ & $3.91$ \\
\bottomrule
\end{tabular}
\caption{The deployed setting is dominated by 8 of its 9 alternatives on both axes at once.}
\end{table}

It is not a latency/accuracy trade-off; it is off the efficient frontier.

The mechanism links this section to Section~\ref{sec:stats}. At $\tau = 3$~s consecutive
outputs are nearly identical ($\rho = 0.962$), yielding $1.2$ independent observations per
minute---which is precisely the coefficient that produces an effective sample size of $25$
from a $20$-minute session. The smoother is simultaneously the largest latency term and the
dominant cause of the autocorrelation that collapses statistical power. Shortening it pays
twice.

\subsection{The floor, and what binds it}

A streaming estimator has no window centroid and no hop quantisation; its only delays are
filter group delay and the smoother. At second order with $\tau = 0.1$~s that is $0.189$~s.

\begin{table}[h]
\centering
\begin{tabular}{lrrr}
\toprule
analysis path & crossfade & end-to-end & vs the $500$~ms result \\
\midrule
deployed, $5.500$~s & $1.00$~s & $6.500$~s & $13.0\times$ \\
retuned, $1.750$~s & $1.00$~s & $2.750$~s & $5.5\times$ \\
floor, $0.189$~s & $1.00$~s & $1.189$~s & $2.4\times$ \\
floor, $0.189$~s & $0.25$~s & $0.439$~s & $0.9\times$ --- inside \\
\bottomrule
\end{tabular}
\caption{The architecture's floor reaches the regime in which alpha neurofeedback still
works.}
\end{table}

Two consequences. ``Analysis is $85\%$ of the budget'' is configuration-dependent, not
architectural: fix the analysis path and the crossfade becomes the binding term, inverting
the decomposition we report above. And the floor is not a trade-off: per-sample
discriminability collapses there ($d = 1.99 \rightarrow 0.61$) but independence rises from
$1.2$ to $41.4$ observations per minute, and information per minute at the floor is
$1.83\times$ the deployed value. The fastest configuration this architecture can reach is
also the more informative one.

Three caveats. Measured \emph{detection} does not fall monotonically with the budget---below
$\tau \approx 0.25$ the estimate is noisy enough that midpoint crossing becomes harder to
time, so the budget is a floor on delay rather than a promise about detection. Every figure
here is measured on an eyes-open/closed manipulation, which is gross compared to
adaptive-versus-sham. And no session has been run at these settings: the frontier is located
and priced on a proxy, not paid for.

\subsection{A finer estimator needs controller work, not a flag}

Adopting the faster estimator is not the two-parameter change it appears to be. Replayed
against a pilot recording, $2$~s${}/{}0.5$~s${}/{}\tau = 0.5$ produces $591$ switches inside
a crossfade---a continuous blend of two independent renders rather than music with
transitions in it, and a larger fraction than the defect that made our first pilot's audio
unusable. The cause is not the estimator: it is that the ladder quantises $z$ with no
hysteresis, so a less-smoothed index crosses rung boundaries constantly.

The fix is a minimum dwell, and its value is structural rather than tuned: a dwell of at
least one crossfade is exactly the condition for no switch arriving before the previous
crossfade completes, so sub-crossfade switches fall to $0$ by construction. The dwell is not
free---the library engine acts within one crossfade, so worst case becomes dwell plus
crossfade---which is why the honest figure for the retuned configuration is $1.8\times$
rather than the $2.4\times$ the estimator alone suggests.

\section{Four statistical properties of closed-loop EEG data}
\label{sec:stats}

\subsection{Autocorrelation inflates evidence by a factor of six}

A $20$-minute session yields $1043$ valid windows at a $1$~s hop, and it is tempting to treat
them as $1043$ observations. Lag-1 autocorrelation is $0.953$, giving an AR(1) effective
sample size of $25.3$. Any test treating windows as independent overstates its evidence by
$\sqrt{1043/25.3} = 6.4\times$---the difference between $p = 0.05$ and $p = 0.4$.

\subsection{Continuous coupling loses power exactly when the intervention works}

Cross-correlating the audio envelope against the neural index over a whole session has a
perverse property: a participant who reaches target and stays there keeps the controller on
one rung, so the audio stops varying in the way the controller drives it and there is almost
nothing left to correlate. Our first pilot returned $r = -0.054$, $p = 0.795$ while sitting
on one rung for $95.9\%$ of the session. The estimator was not failing---it recovers a known
$+6$~s lag from synthetic data. There was simply no controller-driven variation in the window
it was given. This is a bad property for a primary outcome and it is not fixable by improving
the estimator.

\subsection{Event-locked responses are confounded by their own trigger}

Conditioning on rung changes restores power, because it looks only at moments that carry
information. But a rung change happens \emph{because} $z$ moved, and $z$ keeps moving---so a
positive effect is exactly what a closed loop with no therapeutic effect whatsoever would
produce.

We calibrate this rather than warn about it. Against synthetic sessions with known ground
truth, the estimator recovers $+1.503$ on a responder, $-1.467$ on an anti-responder, and
$+0.011$ ($p = 0.573$) on an unrelated session. A fourth session is constructed with no audio
effect at all---$z$ wanders and rung changes fire on boundary crossings, as the real
controller does---and it scores $+0.435$ at $p = 0.010$.

Our pilot measures $+0.412$. The two are indistinguishable. The diagnostics separate them
from a real response and the pilot sits on the wrong side: its pre-window slope is $+0.1292$
z/s against $+0.0573$ for the pure confound and $+0.0012$ for a genuine response, and it is
the only one of the three whose curve peaks at onset and decays---the shape a trailing
excursion makes, and the opposite of a real response, which rises \emph{after} onset. A
positive within-arm event-locked effect is not evidence of an effect. Only the yoked contrast
is interpretable.

\subsection{Dichotomising costs about three times the effect}

Time-in-band is the interpretable clinical quantity and a costly primary outcome: it discards
how far inside or outside the band a participant is and saturates, since someone well outside
the band scores near zero in both arms. Simulated against the measured autocorrelation, it
costs roughly $3\times$ in detectable effect. We register mean $z$ as primary and report
time-in-band descriptively.

\paragraph{Consequence for design.} At $n = 10$ with a crossover design, power to detect the
literature-matched $0.15$~$z$ effect is $38.8\%$; reaching $80\%$ needs $25$ participants per
arm. A study reporting $p > 0.05$ at $38\%$ power has learned nothing, so we register a
feasibility study reporting an interval and two variance components, and no $p$-value for the
primary contrast.

\section{Does the index measure anything?}

Two separable claims: that the electrodes record cortex, and that $\log(\beta/\alpha)$ tracks
arousal.

\paragraph{Cortex.} On an eyes-open/eyes-closed manipulation, alpha rose $2.13\times$ with eye
closure on the temporal pair ($d = 1.55$, $p = 2.8 \times 10^{-23}$). That figure is
threshold-dependent and we report the sensitivity rather than the single number: the recording
used a $150~\mu$V rejection threshold, which removes eyes-open windows preferentially (blinks
occur with the eyes open) and leaves a $1.97$ imbalance; at the pipeline's deployed
$350~\mu$V threshold the conditions are near-balanced and the ratio is $1.90\times$. The
effect is significant at every threshold tested including no rejection at all, so it is not
manufactured by the rejection rule---only its magnitude is.

\paragraph{The channels the study actually uses are weaker, and we report both recordings
rather than substituting one.} The validation above is on TP9/TP10; the index runs on
AF7/AF8. A second recording on AF7/AF8 with good contact ($6.8\%$ rejection) shows a
significant alpha rise---$1.34\times$ geometric, $d = 0.61$, $p = 5 \times 10^{-7}$---but
fails a stricter diagnostic: the eyes-closed rise in \emph{spectral peak prominence}, which
the temporal channels pass clearly ($2.46\times$, $1.91\times$), is $1.09\times$ on AF7 and
$1.13\times$ on AF8, below a $1.2\times$ threshold. Alpha power rises on the forehead; a
distinct alpha peak barely does. We report frontal signal quality as a measured limitation.
This is consistent with Monroy-D'Croz et al., whose frontal alpha asymmetry explained $0.40\%$
of signal variance.

\paragraph{Arousal.} Tested on DEAP, which pairs 32-channel EEG with self-reported affect.
Computed with the same feature extractor the live system uses, the index correlates with
reported arousal at $\rho = +0.303$ ($p = 0.058$) on AF3/AF4, against $+0.082$ with
valence---the right discriminant profile, at a sample size where no single montage reaches
significance. The direction holds across 7 of 7 independent montages (sign test $p = 0.016$),
which is the result; no single correlation is.

\section{Reproducibility as an engineering problem}

A preprint's numbers are copied from a terminal into prose once, and then the code moves.
Every headline number in this paper is regenerated from artefacts on disk by a single command,
which checks it against the value asserted here: $44$ claims, all reproducing. The practice
caught real drift, including one instance where our own motivating example (a manuscript
saying $1.14\times$ while the repository said $1.05\times$) was live in the codebase.

Three failure modes were worth the engineering, and we report them because they generalise to
any pre-registered study with a live codebase.

\emph{Recordings must be pinned by name, not selected as ``the newest''.} A second pilot
silently re-based a dozen claims, five of seven figures and eight tests onto itself---including
a claim that documented a \emph{historical} defect, which would have been overwritten with the
post-fix number and the finding quietly deleted.

\emph{Synthetic data must be unable to reach the manuscript.} A thirty-second demo run wrote a
session that recomputed six claims from a signal generator. The marker distinguishing
synthetic from real already existed; nothing read it.

\emph{The freeze mechanism must not force its own violation.} Our pre-registration is
hash-checked on every test run, but its deviations section is \emph{inside} the hashed file,
so logging a deviation fails the check. The obvious repair---updating the stored hash---destroys
the guarantee permanently. A guard that makes disabling itself the reasonable next step is
worse than no guard; the log lives outside the hash.

Two data-acquisition defects are also worth reporting, because both were invisible to every
check we had. The streaming layer delivered each packet two to four times with timestamps
stepping backwards between copies, which downstream code---counting samples rather than
time---read as the signal jumping backwards every $12$ samples and reported as electrode
failure on a headset that was fine. And raw timestamps were stored at a precision that
depended on how long the laptop had been running: $8$~ms after a fresh boot, $0.25$~s after
$27$ days up.

\section{What this work does not establish}

No efficacy claim of any kind is made or supported. The evidence base is two
self-administered unblinded sessions, zero adaptive/sham pairs and zero participants.

\paragraph{The measured latency is far outside the regime where the delay literature finds
effects.} Our deployed budget is $13\times$ the delay at which sustained alpha change had
already disappeared in Belinskaia et al., and the retuned configuration is still $5.5\times$.
We do not think this invalidates the approach, but the argument must be made rather than
assumed: operant neurofeedback asks the participant to perceive a contingency, which is where
the delay constraint is derived, whereas the proposed mechanism here is the iso-principle,
whose evidence comes from fixed sequences with no contingency at all---Starcke et
al.~\cite{starcke} randomised $107$ participants to four pre-set sequences after a sadness
induction and found the iso sequence produced higher positive affect and valence
($\eta^2 = 0.08$, $0.09$) with purely passive listening. That argument has a sharp edge we
state rather than hide: if matching-then-leading works with a fixed playlist, the burden on a
closed-loop system is not to beat silence but to beat a fixed playlist---which is exactly what
a yoked sham tests.

\paragraph{A controller can be well-behaved and still produce no manipulation.} Our second
pilot ran for twenty minutes with zero prompt changes: at $1.0$~SD per rung, a participant who
sits at or below target maps to a single rung for the whole session. A yoked sham of such a
session is acoustically identical to the adaptive arm, so the contrast has nothing to test. We
widened the ladder to nine rungs at $0.5$~SD, which restores variation ($33$ changes across
$5$ prompts on one pilot, $26$ across $3$ on the other)---but this was discovered only by
running the system on a person and counting what it played. We suggest reporting
prompt-change counts and distinct-prompt counts as standard for any adaptive audio
intervention; they are cheap, and their absence hides a degenerate manipulation behind a
functioning system.

\section{Conclusion}

We measured what a closed-loop EEG music system actually costs in time, and found the dominant
term in the analysis path rather than the model---then found that this decomposition is a
property of the configuration, not the architecture. The deployed estimator is dominated by 8
of 9 alternatives on both latency and information rate; the architecture's floor is $0.189$~s
of analysis and $0.439$~s end-to-end, inside the regime where alpha neurofeedback still works,
and more informative than the configuration we shipped.

Alongside, four statistical properties of closed-loop EEG data that shaped every design
decision we made: autocorrelation inflating evidence $6.4\times$, continuous coupling failing
when the intervention succeeds, event-locked effects indistinguishable from their own trigger
confound, and dichotomisation costing $3\times$ in detectable effect.

None of it says the intervention works. All of it is what we would have wanted to read before
building the system.

\paragraph{Reproducibility.} $44$ of the numbers in this paper are regenerated from artefacts
on disk by \texttt{python scripts/verify\_claims.py}, which checks each against the value
asserted here and fails if any has moved. The synthetic-ground-truth figures in
Section~\ref{sec:stats} come from \texttt{validate\_event\_locked.py}, the coupling recovery
from \texttt{validate\_coupling.py}, and the threshold sweep from
\texttt{alpha\_sensitivity.py}. Effect sizes quoted from the cited literature are not ours to
regenerate.

\begin{thebibliography}{99}

\bibitem{mindmelody} Y. Zhang, Y. Sun, H. Gu and Z. Jin.
\emph{MindMelody: A Closed-Loop EEG-Driven System for Personalized Music Intervention.}
arXiv:2605.01235, 2026.

\bibitem{mindtomusic} Ran et al.
\emph{Mind to Music: An EEG Signal-Driven Real-Time Emotional Music Generation System.}
International Journal of Intelligent Systems, 9618884, 2024. doi:10.1155/int/9618884

\bibitem{ehrlich} S. K. Ehrlich, K. R. Agres, C. Guan and G. Cheng.
\emph{A closed-loop, music-based brain-computer interface for emotion mediation.}
PLOS ONE, 2019. doi:10.1371/journal.pone.0213516

\bibitem{dashagres} A. Dash and K. R. Agres.
\emph{AI-Based Affective Music Generation Systems: A Review of Methods, and Challenges.}
arXiv:2301.06890, 2023.

\bibitem{neurophone} Velasquez-Pena et al.
\emph{Neurophone: A Brain-Computer Music Interface for Emotional Neurofeedback.}
IEEE International Symposium on the Internet of Sounds, 2025.

\bibitem{minimalist} P. A. Monroy-D'Croz, R. Ramirez-Melendez and J. Cespedes-Guevara.
\emph{A Minimalist Brain-Computer Musical Interface for Real-Time Emotion-Driven
Sonification.} arXiv:2606.01473, 2026.

\bibitem{belinskaia} A. Belinskaia, N. Smetanin, M. Lebedev and A. Ossadtchi.
\emph{Short-delay neurofeedback facilitates training of the parietal alpha rhythm.}
Journal of Neural Engineering, 2020. doi:10.1088/1741-2552/abc8d7

\bibitem{asai} T. Asai, T. Hamamoto, S. Kashihara and H. Imamizu.
\emph{Real-Time Detection and Feedback of Canonical Electroencephalogram Microstates.}
Frontiers in Systems Neuroscience, 2022. doi:10.3389/fnsys.2022.786200

\bibitem{gesture2music} R. Jeyaraj, B. Subramanian, K. Gangadharan and A. Paul.
\emph{Gesture2Music: A Low-Latency Real-Time Framework for Continuous Gesture-Driven Music
Generation.} arXiv:2511.00793.

\bibitem{altshuler} I. M. Altshuler.
\emph{The past, present and future of psychiatric music therapy.}
American Journal of Psychiatry, 1948.

\bibitem{starcke} K. Starcke, J. Mayr and R. von Georgi.
\emph{Emotion Modulation through Music after Sadness Induction---The Iso Principle in a
Controlled Experimental Study.} International Journal of Environmental Research and Public
Health, 2021. PMC8656869.

\end{thebibliography}

\end{document}
